\documentclass{article}

\usepackage{a4}
\usepackage{url}
\usepackage{amsmath}
\usepackage{setspace}

\textwidth 6in
\oddsidemargin 0.15in
\topmargin -0.5in
\textheight 10in
\doublespacing

\begin{document}

\pagestyle{empty}
\thispagestyle{empty}

\begin{center}
    {\LARGE \bf \noindent Urdu Text Detection In Natural Scene Images}
\end{center}

\vspace{0.2in}
{\bf \noindent Abstract}\\
\vspace{-0.1in}

\noindent  This thesis embarks on a comprehensive exploration of text detection within complex natural scenes, with a focal point on the Urdu language. Urdu, a widely spoken South Asian language, introduces unique complexities due to its cursive script and diverse writing styles. The primary research objectives revolve around developing robust Urdu text detection algorithms, enhancing the detection of multi-oriented and arbitrary-shaped text, and facilitating their integration into real-world applications. The contributions of this work encompass dataset generation to enrich the research landscape and the introduction of an innovative architecture that combines the Boundary Proposal Model and the Adaptive Deformation Model.

The experimental journey involves the evaluation of three models - YOLO, Faster R-CNN, and a two-fold model, the Boundary Proposal and Adaptive Deformation Model. While each model demonstrates substantial text detection capabilities, the Boundary Proposal and Adaptive Deformation Model excels by effectively identifying arbitrary, multi-oriented, and complex-shaped text, even within intricate backgrounds. Acknowledging the model's limitations, particularly in scenarios with severe blurring and occlusions, this research transcends its focus on a single language or script. It extends its scope to embrace the realm of multilingual text detection and encourages cross-domain collaboration to foster innovation in text detection technologies. The imperative driving this work is to bridge technological gaps, advancing text detection to align with the demands of emerging technologies such as autonomous devices. This research opens new horizons for practical applications, including autonomous systems, translation services for cross-cultural interactions, and enhanced text accessibility across diverse scenarios.

\end{document}

